\section{Related literature}
\label{sec:rewrite-literature}

This paper addresses the connected questions of growth, income distribution, and transition dynamics in the research agenda proposed by \citet{brynjolfssonkorinekagrawal2025}. My closest antecedents are models in which AI improves both production and research. \citet{davidsonetal2026} and \citet{jonestonetti2026} study feedback between technological progress, research, and output. Technological progress is endogenous in both papers, but saving and research allocation shares are taken as given. Another close antecedent is \citet{erdiletal2025gate}, who model how investment in compute, training, and technology expands both the range of tasks AI can perform and the services it supplies within that range. Whereas they solve a social planner's problem, I study a market equilibrium with a monopolistic AI developer. I likewise distinguish resources used to deploy AI from those used to improve it. I focus on how RSI increases the services obtained per unit of compute, without explicitly modeling the automation of additional tasks.
%They let a social planner choose investment and compute allocation. I, in contrast, study a decentralized economy in which households choose saving, competitive firms choose production inputs, and a monopolistic AI developer chooses research investment to maximize discounted net profit. The resources devoted to improving AI therefore respond to private profit opportunities and compete with consumption and physical-capital accumulation.

The developer's research incentives build on \citet{romer1990}, where
expected monopoly profits motivate investment in nonrival knowledge.
\citet{liuwan2026} study the interaction between a general-purpose AI
supplier and developers of specialized applications. In their baseline
model, a higher markup on general-purpose AI services can strengthen the
supplier's research incentive while raising research costs for application
developers. I model a single AI developer, without distinguishing
general-purpose AI from specialized applications.

The distinction between temporary and permanent growth gains relates to
the production bottlenecks studied by \citet{aghionjonesjones2019} and
\citet{jonestonetti2026}: tasks that remain dependent on human labor can
constrain growth despite automation elsewhere.
%I hold the aggregate production weights fixed and characterize how substitution between labor and AI services, together with an upper bound on AI efficiency, determines the long-run regime.
In my model, RSI need
not eliminate the labor bottleneck. Conversely, when AI and labor are
substitutes and AI is already sufficiently efficient, permanently faster
growth can persist without RSI, since capital and compute can
continue expanding together. This connects to \citet{restrepo2025agi},
where output can keep growing with the exogenous expansion of compute
once bottleneck tasks have been automated.
%When AI and labor are substitutes, I identify an efficiency threshold above which AI-dominated equilibria exist even though AI efficiency approaches an upper bound.

\citet{acemoglurestrepo2018} show that automation can raise the long-run
wage level while reducing labor's income share. \citet{raymookherjee2022}
show that capital accumulation can sustain rising wages even as labor's
share vanishes. In \citet{hemousolsen2022}, low-skill wages grow more
slowly than output, while high-skill wages keep pace with it. In my
AI-dominated regime, wages grow permanently faster than in the
labor-bottleneck regime, but more slowly than output per worker.
The substitution elasticity determines how much of the additional
output-per-worker growth translates into faster wage growth.
%I characterize these outcomes under private investment in RSI. In the AI-dominated equilibria with an upper bound on AI efficiency, wage growth permanently exceeds the exogenous rate of labor-augmenting technological progress, but output per worker grows even faster. An exact relationship links this wage-growth premium to the substitution elasticity and the rate at which labor's share declines. Separating AI revenue from compute expenditure and net profit also shows how the industry's growing share of output is distributed between its resource costs and its owners' income.

\citet{korinekstiglitz2019} distinguish the gains earned by innovators
from the redistribution caused by changes in wages and other factor
prices. \citet{korinekng2019} show how digital innovation can increase
monopoly profits as a share of income while reducing the shares paid
to labor and physical capital. I likewise distinguish the AI developer's
net profit from the remuneration of physical capital, and examine how
these incomes evolve alongside wages as AI efficiency improves and
capital accumulates.

\citet{mollrachelrestrepo2022} show that automation can raise returns to
wealth and increase inequality across households. In my AI-dominated
regime, the return to physical capital and output-per-worker growth both
exceed their labor-bottleneck values. I study the division of income between labor, physical
capital, and AI-industry profits, but not inequality across households
with different asset holdings.
